\chapter{Metodología}
\label{ch:metodologia}

Este capítulo integra el diseño metodológico completo: conjuntos de datos, protocolos de simulación de canales faltantes, construcción de grafos, algoritmos de interpolación, métricas de evaluación, motor de optimización bayesiana, arquitectura del sistema experimental, trazabilidad de artefactos, diseño experimental y protocolo de validación comparativa. La Figura~\ref{fig:methodology_flow} resume el flujo metodológico completo.

\begin{figure}[htbp]
\centering
\includegraphics[width=0.80\textwidth]{ch3_methodology_flow}
\caption[Flujo metodológico de evaluación]{Flujo metodológico de evaluación con control pareado. Cada método se aplica sobre los mismos conjuntos de datos, máscaras de pérdida y semillas; luego se calculan métricas de reconstrucción y costo antes de comparar los métodos por escenario pareado.}
\label{fig:methodology_flow}
\end{figure}
% ======================================================================
\section{Conjuntos de Datos}
% ======================================================================

La evaluación se realizó sobre tres bases de datos EEG públicas que cubren un espectro diverso de densidades de electrodos, paradigmas y poblaciones (Tabla~\ref{tab:datasets}).

\begin{table}[htbp]
\centering
\caption{Características de los tres conjuntos de datos experimentales.}
\label{tab:datasets}
\small
\begin{tabular}{@{}lccccc@{}}
\toprule
Dataset & Canales & Sistema & Frec. (Hz) & Sujetos & Paradigma \\
\midrule
PhysioNet & 64 & 10--10 & 160 & 109 & Reposo \\
BCI IV 2a & 22 & 10--20 & 250 & 9 & Imag. motora \\
MNE Sample & 60 & 10--20 ext. & 600 & 1 & Auditivo \\
\bottomrule
\end{tabular}
\end{table}

\textbf{PhysioNet} \cite{goldberger2000physiobank,schalk2004bci2000} contiene 64 canales (10--10) a 160~Hz de 109 sujetos. Se usaron segmentos de reposo con ojos abiertos (60~s por sujeto). \textbf{BCI Competition IV 2a} \cite{brunner2008bci} aporta 22 canales (10--20) a 250~Hz de 9 sujetos en tareas de imaginación motora de cuatro clases. Se extrajeron 2592 ensayos de 2~s. \textbf{MNE Sample} \cite{gramfort2013meg} corresponde a un registro de un sujeto; de ese registro se seleccionaron 60 canales EEG a 600~Hz y 72 ensayos de 0.7~s alrededor de estímulos auditivos.

Estos tres conjuntos cubren el espacio de variabilidad real: desde 22 canales (clínica ambulatoria) hasta 64 (investigación), con frecuencias de 160 a 600~Hz y paradigmas de reposo, imaginación motora y potenciales evocados.
% ======================================================================
\section{Protocolos de Simulación de Canales Faltantes}
% ======================================================================

La evaluación de métodos de interpolación requiere conocer la señal verdadera (\textit{ground truth}) para cuantificar el error de reconstrucción. La estrategia adoptada consiste en simular la pérdida sobre canales que sí están presentes y cuya señal real se conoce: se selecciona un subconjunto de canales, se oculta su señal (máscara de falta), se aplica el método de interpolación para estimar dicha señal, y se compara la estimación con la señal original oculta. Esta estrategia ---estándar en la literatura de imputación de señales fisiológicas--- permite calcular métricas de error exactas sobre datos reales sin requerir registros simultáneos de referencia.

Se definieron dos modos de pérdida de canales. El \textbf{modo aleatorio} selecciona un porcentaje $p$ de canales al azar, modelando fallas esporádicas (secado del gel, desprendimiento de electrodos). El \textbf{modo contiguo} selecciona un canal semilla y marca como faltantes los $p$ canales más cercanos en distancia euclidiana, modelando fallas catastróficas localizadas (saturación de amplificador).

Para cada modo se evaluaron cuatro severidades ($p \in \{10\%, 20\%, 30\%, 40\%\}$) con 5 semillas aleatorias, generando $3 \times 2 \times 4 \times 5 = 120$ escenarios. La semilla determina qué canales específicos se marcan como faltantes y se preserva para garantizar que todos los métodos se evalúen sobre exactamente las mismas máscaras.

% ======================================================================
\section{Construcción de Grafos}
% ======================================================================
\label{sec:graph_construction_methods}

Los métodos basados en GSP requieren un grafo subyacente que define la topología de conectividad entre electrodos. La construcción del grafo es una decisión metodológica independiente del algoritmo de interpolación: un mismo método puede operar sobre distintos grafos, y un mismo grafo puede alimentar distintos algoritmos.

Se implementaron diez estrategias de construcción de grafos, organizadas según la fuente de información que utilizan (Tabla~\ref{tab:graph_methods}). La nomenclatura entre paréntesis corresponde al identificador usado en el código del repositorio.

\begin{table}[htbp]
\centering
\caption{Estrategias de construcción de grafos. La columna «Fase» indica si el grafo se usó en la fase exploratoria (E), comparativa (C) o ambas (E+C).}
\label{tab:graph_methods}
\small
\begin{tabular}{@{}llp{4.8cm}l@{}}
\toprule
Estrategia & Tipo & Descripción & Fase \\
\midrule
knn & Geométrico & k-NN no ponderado ($k$ vecinos, pesos binarios) & E \\
knng & Geométrico & k-NN con kernel gaussiano ($k$, $\sigma$) & E+C \\
vknng & Geométrico & k-NN variable según densidad local ($\alpha$) & E \\
epsilon\_ball & Geométrico & $\varepsilon$-ball (todos los pares con $d<\varepsilon$) & E \\
gaussian & Geométrico & Distancia completa con kernel gaussiano ($\sigma$) & E+C \\
inv\_distance & Geométrico & Conectividad total con $W_{ij}=1/d_{ij}$ & E \\
correlation & Correlación & $W_{ij}=|\rho_{ij}|$ (solo fase exploratoria) & E \\
kalofolias & Data-driven & Optimización de log-grados + regularización Frobenius & E \\
nnk & Data-driven & Regresión de kernel no-negativa (NNLS local) & E \\
aew & Data-driven & \textit{Adaptive Edge Weighting} & E \\
\bottomrule
\end{tabular}
\end{table}

Los grafos geométricos operan exclusivamente sobre las posiciones físicas de los electrodos, lo que los hace inmunes a canales faltantes y deterministas. Las posiciones se obtuvieron desde los montajes disponibles en MNE-Python: PhysioNet se estandarizó con el montaje \code{standard\_1005}; BCI IV 2a se mapeó a los 22 nombres de canal del montaje 10--20 extendido; y MNE Sample usó las coordenadas incluidas en el archivo \code{sample\_audvis\_raw.fif}. Cuando un canal no tenía coordenada explícita, el sistema lo marcó como caso a revisar y evitó usarlo para conclusiones fuertes. Los grafos de correlación funcional se usaron únicamente en la fase de exploración inicial (\textit{screening}), donde era necesario evaluar un gran número de combinaciones método$\times$grafo; se descartaron para la fase comparativa porque la correlación se estima sobre las mismas señales que luego se reconstruyen, introduciendo circularidad metodológica. Los grafos data-driven se evaluaron en la fase exploratoria pero no pasaron a la fase comparativa por su mayor costo computacional y menor reproducibilidad determinista. Los grafos knng y gaussian fueron los únicos que avanzaron a la fase comparativa, por su balance entre expresividad y reproducibilidad.
% ======================================================================
\section{Métodos de Interpolación}
% ======================================================================
\label{sec:interpolation_methods}

Se implementaron dieciocho métodos de interpolación organizados en cuatro familias (Tabla~\ref{tab:methods}). La nomenclatura de familias es consistente a lo largo de todo el documento.

\begin{table}[htbp]
\centering
\caption{Métodos de interpolación implementados y evaluados.}
\label{tab:methods}
\small
\begin{tabular}{@{}lp{6.5cm}l@{}}
\toprule
Familia & Métodos & Fase \\
\midrule
Geométrica & Vecino más cercano, Distancia inversa (IDW), RBF, Spherical splines, Lineal & E \\
Estadística & ICA (FastICA), ICA-MNE (Picard) & E \\
GSP & Tikhonov en grafos, Banda limitada (BGSRP), Suavizado en grafos, TV en grafos, GTT (graph-time Tikhonov) & E \\
TRSS & Regularización espacio-temporal de Sobolev & E+C \\
Referencia & MNE-Python (defecto), MNE-Python (Optuna) & E+C \\
\bottomrule
\end{tabular}
\end{table}

Es metodológicamente crucial distinguir entre la implementación propia de \textit{spherical splines} y la de MNE-Python, que incorpora mejoras de ingeniería documentadas en el Capítulo~\ref{ch:marco_teorico}. Los métodos GSP (Tikhonov, BGSRP, TV, GTT) y TRSS se evaluaron sobre múltiples grafos de la Sección~\ref{sec:graph_construction_methods}, generando variantes por método según el grafo subyacente. Los métodos geométricos y estadísticos se evaluaron en la fase exploratoria; solo TRSS y MNE-Python (en sus variantes por defecto y optimizada) avanzaron a la fase comparativa.

% ======================================================================
\section{Métricas de Evaluación}
% ======================================================================
\label{sec:metricas}

Se emplearon seis métricas primarias organizadas en tres categorías, más el tiempo de cómputo como métrica complementaria (Tabla~\ref{tab:metrics}). Ninguna métrica individual captura todos los aspectos relevantes; por ello se reportan de forma independiente en los resultados.

\begin{table}[htbp]
\centering
\caption{Métricas de evaluación. La notación $\mathcal{F}$ denota el conjunto de canales faltantes, $x_i$ la señal real y $\hat{x}_i$ la reconstrucción.}
\label{tab:metrics}
\footnotesize
\begin{tabular}{@{}p{2.0cm}p{6.7cm}cc@{}}
\toprule
Categoría & Métrica & Símbolo & Dirección \\
\midrule
\multirow{3}{*}{Amplitud} 
& Error absoluto medio: $\frac{1}{|\mathcal{F}|}\sum_{i\in\mathcal{F}}|x_i-\hat{x}_i|$ 
& MAE & $\downarrow$ \\
& Raíz del error cuadrático medio: $\sqrt{\frac{1}{|\mathcal{F}|}\sum_{i\in\mathcal{F}}(x_i-\hat{x}_i)^2}$ 
& RMSE & $\downarrow$ \\
& Relación señal-ruido: $10\log_{10}(\|x_{\mathcal{F}}\|^2/\|x_{\mathcal{F}}-\hat{x}_{\mathcal{F}}\|^2)$ 
& SNR & $\uparrow$ \\
\midrule
Morfológica 
& \textit{Dynamic Time Warping}: $\mathrm{DTW}(x,\hat{x})$ 
& DTW & $\downarrow$ \\
\midrule
\multirow{2}{*}{Espectral} 
& Distancia Espectral Logarítmica: $\sqrt{\frac{1}{B}\sum_b (\log S_x(b)-\log S_{\hat{x}}(b))^2}$ 
& LSD & $\downarrow$ \\
& Coherencia: $|S_{x\hat{x}}(f)|^2/(S_{xx}(f)S_{\hat{x}\hat{x}}(f))$ 
& Coh. & $\uparrow$ \\
\midrule
Temporal & Tiempo de cómputo por caso (s) & $t$ & $\downarrow$ \\
\bottomrule
\end{tabular}
\end{table}

Las métricas de \textbf{amplitud} cuantifican el error muestra a muestra sobre los canales faltantes. El \textbf{DTW} permite deformaciones no lineales del eje temporal, relevante para segmentos con transitorios. Las métricas \textbf{espectrales} (LSD, Coherencia) cuantifican la distorsión del espectro de forma independiente del error de amplitud. El \textbf{tiempo de cómputo} se reporta como dimensión complementaria porque la mejora en precisión de reconstrucción debe sopesarse contra el costo computacional adicional.

% ======================================================================
\section{Optimización Bayesiana con Optuna}
% ======================================================================
\label{sec:optuna}

Para garantizar una comparación justa, todos los algoritmos fueron optimizados mediante Optuna \cite{akiba2019optuna} con muestreo TPE (\textit{Tree-structured Parzen Estimator}). El TPE modela las distribuciones de configuraciones con buen y mal desempeño y guía la búsqueda hacia regiones prometedoras, siendo más eficiente que una búsqueda aleatoria o de grilla para espacios de alta dimensionalidad.

Los espacios de búsqueda se definieron por familia (Tabla~\ref{tab:hyperparams}). Se ejecutaron 30 intentos por método de grafo, 10 intentos con NSGA-II multiobjetivo para \textit{spherical splines}, y barridos de grilla fina para TRSS. En total se exploraron más de quince mil configuraciones paramétricas distintas.

\begin{table}[htbp]
\centering
\caption{Espacios de búsqueda de hiperparámetros.}
\label{tab:hyperparams}
\small
\begin{tabular}{@{}l c c l@{}}
\toprule
Familia & Parámetros & Escala & Estrategia \\
\midrule
GSP (grafo) & $\sigma \in [0.01, 1.0]$, $k \in \{3,5,7,10,15,20\}$ & Lineal & TPE (30 intentos) \\
TRSS & $\alpha \in [10^{-4}, 10^2]$, $\beta \in [10^{-3}, 10^3]$ & Log & Grilla fina \\
Splines propios & $m \in \{2,3,4\}$, $\lambda \in [10^{-8}, 10^{-2}]$ & Mixta & NSGA-II (10 intentos) \\
MNE (Optuna) & $m \in \{2,3,4\}$, $\lambda \in [10^{-8}, 10^{-2}]$ & Mixta & NSGA-II (10 intentos) \\
\bottomrule
\end{tabular}
\end{table}

% ======================================================================
\section{Arquitectura del Sistema Experimental}
% ======================================================================
\label{sec:architecture}

El sistema se organiza en cinco módulos con interfaces explícitas que operan en cadena (Figura~\ref{fig:architecture_blocks}). El módulo de \textbf{datos} expone una interfaz común para los tres conjuntos. El módulo de \textbf{simulación de daños} genera máscaras binarias con semilla trazable. El módulo de \textbf{métodos} combina la construcción de grafos —diez estrategias que producen $\mathbf{W}$ y $\mathbf{L}$— con las dieciocho implementaciones de interpolación bajo una interfaz unificada. El módulo de \textbf{evaluación} calcula las métricas primarias y complementarias, y el módulo de \textbf{artefactos/optimización} orquesta Optuna en la fase exploratoria y registra resultados trazables para la fase comparativa.

\begin{figure}[H]
\centering
\includegraphics[width=0.98\textwidth]{ch3_architecture_blocks}
\caption[Diagrama de bloques de la arquitectura]{Diagrama de bloques de la arquitectura del sistema experimental. Cinco módulos con interfaces explícitas operan en cadena. La figura distingue fase exploratoria y fase comparativa, indicando que los resultados comparativos no reabren la búsqueda de hiperparámetros.}
\label{fig:architecture_blocks}
\end{figure}
\subsection{Trazabilidad y Gestión de Artefactos}

En experimentos computacionales a gran escala (18 métodos, 120 escenarios, más de 15\,000 configuraciones), la gestión sistemática de artefactos resuelve tres problemas concretos. Primero, \textbf{auditoría}: cada número reportado debe poder rastrearse hasta el archivo que lo origina; sin trazabilidad es imposible distinguir un hallazgo real de un error de transcripción. Segundo, \textbf{control de contaminación}: las fases exploratoria y comparativa deben mantenerse físicamente separadas para que resultados exploratorios no influyan inadvertidamente en decisiones posteriores. Tercero, \textbf{reproducibilidad}: un tercero con acceso al repositorio (\url{https://github.com/csaldivia/Tesis-GSP-EEG}) puede re-ejecutar cualquier subconjunto del experimento.

Cada ejecución recibe un identificador único (SHA-256 de la tupla completa de parámetros) y preserva: configuración inmutable (YAML), métricas agregadas (JSON), señales reconstruidas en formato NPZ (NumPy comprimido), máscara de daño, y tiempos de ejecución (Tabla~\ref{tab:artifact_types}). El código está versionado con Git y cada conjunto de resultados está asociado a un \textit{commit hash} específico. La Figura~\ref{fig:traceability_pipeline} ilustra la cadena de trazabilidad desde el texto hasta el script generador.

\begin{figure}[htbp]
\centering
\includegraphics[width=0.96\textwidth]{ch4_traceability_pipeline}
\caption[Cadena de trazabilidad experimental]{Cadena de trazabilidad experimental. Cada afirmación cuantitativa se vincula con la tabla o figura correspondiente, el archivo de datos asociado y el script que generó el resultado.}
\label{fig:traceability_pipeline}
\end{figure}
\begin{table}[htbp]
\centering
\caption{Artefactos generados por el sistema experimental.}
\label{tab:artifact_types}
\small
\begin{tabular}{@{}l l l@{}}
\toprule
Artefacto & Formato & Función \\
\midrule
Configuración & YAML & Registro inmutable de parámetros \\
Métricas & JSON & Resultados de las métricas requeridas \\
Señal reconstruida & NPZ (NumPy) & Salida verificable bit a bit \\
Máscara de daño & NPZ (NumPy) & Trazabilidad de canales faltantes \\
Tiempos & JSON & Perfil de eficiencia computacional \\
\bottomrule
\end{tabular}
\end{table}

% ======================================================================
\section{Diseño Experimental}
% ======================================================================

El diseño experimental se resume en la Tabla~\ref{tab:research_questions_experiments}, que vincula cada pregunta de investigación con el experimento diseñado para responderla y la evidencia esperada. La lógica central es pareada: dos métodos se comparan sobre la misma señal, la misma máscara de canales ocultos y la misma semilla, evitando que diferencias del caso experimental contaminen la comparación.

La ejecución no consistió en un único benchmark final. Primero se realizó un cribado amplio de combinaciones método--grafo para observar qué familias eran competitivas en datos sintéticos, proxies EEG y registros reales. Luego se redujo el conjunto de candidatos considerando desempeño, costo de ejecución, reproducibilidad y riesgo de circularidad en la construcción del grafo. En una tercera fase se optimizaron hiperparámetros mediante Optuna o grillas controladas y se congelaron las variantes que pasarían a la comparación final. El Capítulo~\ref{ch:resultados} reporta esta trayectoria experimental antes de presentar los resultados finales TRSS--MNE.
\begin{table}[htbp]
\centering
\caption{Relación entre preguntas de investigación, experimentos y evidencia esperada.}
\label{tab:research_questions_experiments}
\small
\begin{tabular}{@{}p{3.2cm} p{3.0cm} p{5.5cm}@{}}
\toprule
Pregunta & Experimento & Evidencia esperada \\
\midrule
P1: ¿Cuándo mejora TRSS? & Comparación global de familias & Reducción de error y ajuste temporal en pérdida agrupada o severa \\
P1: ¿Aporta el término temporal? & Descomposición $\alpha$/$\beta$ & Diferencia observable entre TRSS completo y TRSS sin temporalidad \\
P1: ¿Preserva morfología temporal? & Series temporales y DTW & Menor distorsión temporal en escenarios representativos \\
P2: ¿Cuándo conviene MNE? & Comparación TRSS vs.\ MNE & MNE gana o empata en pérdida baja, señales suaves o restricciones de tiempo \\
P2: ¿Dónde está el punto de cruce? & Curvas por patrón y severidad & Frontera de uso dependiente de severidad, patrón y métrica \\
\bottomrule
\end{tabular}
\end{table}

Las \textbf{variables independientes} son: método de interpolación, conjunto de datos, modo de pérdida, severidad y semilla (5 réplicas). Las \textbf{variables dependientes} son las métricas de la Sección~\ref{sec:metricas}. Las \textbf{variables controladas} incluyen: versiones exactas de todas las bibliotecas, semillas globales, preprocesamiento idéntico (filtro Butterworth 0.5--45~Hz, referencia promedio común, segmentación uniforme), y hardware de ejecución (CPU Intel Core i7-13700H, 32 GB RAM, sin aceleración GPU).

% ======================================================================
\section{Protocolo de Validación Comparativa}
\label{sec:confirmatory}

La validación comparativa es una fase independiente que responde a una necesidad metodológica concreta: los resultados de la fase exploratoria (barridos de 15\,000 configuraciones) no deben usarse como evidencia principal porque el riesgo de sobreajuste es alto. Un método que obtuvo buenos resultados porque sus hiperparámetros se ajustaron extensivamente sobre los datos de prueba no demuestra superioridad general; demuestra capacidad de ajuste al protocolo exploratorio. La fase comparativa reduce este riesgo al congelar los hiperparámetros de TRSS tras la optimización exploratoria y evaluarlo sobre una partición independiente, contra MNE-Python con parámetros por defecto.

La fase se diseñó con 100 estratos definidos por régimen de señal, patrón de pérdida y severidad, con tres semillas por estrato (300 casos pareados). Los regímenes incluyen señales reales y controles sintéticos suave/rugoso. El diseño es de medidas repetidas: cada caso comparte la misma señal y máscara, por lo que las diferencias se atribuyen al método bajo las condiciones del protocolo. El análisis reporta diferencias pareadas, medianas por estrato, tasas de victoria e intervalos de incertidumbre vía bootstrap jerárquico (1000 remuestreos), sin convertir la decisión metodológica en una prueba de hipótesis nula.

% ======================================================================
\section{Control de Calidad}
% ======================================================================

Antes de cualquier análisis se ejecutaron verificaciones automáticas: todas las ejecuciones contienen las métricas requeridas, ningún identificador duplicado con resultados diferentes, las máscaras coinciden con la semilla, y se excluyen ejecuciones con valores no numéricos, tiempo excesivo ($>$10~min) o reconstrucción idéntica a la entrada.

% ======================================================================
\section{Síntesis del Diseño}
% ======================================================================

El diseño experimental de esta tesis se distingue de la literatura previa en cuatro aspectos: todos los métodos fueron optimizados con el mismo motor bayesiano; se incorporó una comparación independiente contra MNE-Python; la evaluación abarca métricas de amplitud, morfología y espectro; y los escenarios incluyen pérdida aleatoria y contigua en un rango de severidades clínicamente relevante. La arquitectura modular, la separación estricta de fases y la trazabilidad completa de artefactos garantizan que las conclusiones sean auditables y reproducibles.

La consecuencia práctica de este diseño es que el Capítulo~\ref{ch:resultados} debe leerse en dos niveles. El primer nivel es experimental: muestra cómo el repositorio pasó de un espacio amplio de grafos y métodos a un subconjunto final defendible. El segundo nivel es comparativo: evalúa ese subconjunto bajo casos pareados y métricas comunes. Esta distinción evita confundir resultados exploratorios con evidencia final y permite justificar por qué la tesis se concentra en TRSS y MNE-Python.

Además, la separación de fases hace explícitas las decisiones negativas del estudio. No todos los métodos implementados avanzaron a la comparación final; algunos fueron descartados por costo, otros por circularidad potencial en la construcción del grafo, y otros por no mejorar respecto de referencias más simples. Documentar esas reducciones es parte del resultado experimental, porque muestra que la comparación final no fue seleccionada arbitrariamente sino derivada de evidencia intermedia trazable.